\section{The Stationary Reader}

A gauge path begins as a finite itinerary of distinguishable readings. It
does not require a completed space of possible histories. It requires only a
ledger of neighboring commitments: this reading can follow that one, this
slip is permitted, this endpoint is held fixed, this residue is still
unspent. The path is finite because the grammar has not earned any other
kind of record. It may name a direction of refinement, but it may not
collapse the direction into an already completed continuum.

The word gauge matters because the path is not merely a list. A list can
forget why its entries can be compared. A gauge path carries the rule by
which one entry is read against the next. It tells the grammar which local
changes are admissible, which labels are only names for the same reading,
and which changes would cross a boundary the current record cannot support.
The gauge therefore quotients away unearned presentation while preserving
the distinctions that still have measurement content.

An action is the finite cost of such a path. It is not an invisible
substance attached to the path, and it is not a completed integral waiting
above the finite record. It is the ledger's way of accumulating the cost of
admissible comparisons. If two adjacent readings agree in the required
sense, the cost reads as balanced there. If they fail to agree, the failure
appears as residue. The action is useful because it lets the grammar ask a
single question of many local comparisons: what remains after the permitted
balances have been paid?

The stationary reader is the rule that asks this question under variation.
A variation is not an arbitrary fantasy path. It is a permitted slip inside
the gauge. The endpoints that have been committed remain committed. The
labels that have been quotiented remain quotiented. Only the degrees of
freedom still open to the grammar may move. This restriction is not a
technical convenience. It is what prevents the reader from using a
variation to erase the very record it is supposed to test.

When a variation is applied, the action changes. The change can be read in
orders. The first order is the linear remainder, the part that would be
visible to an infinitesimal push in any admissible direction. But the
grammar does not need an infinitesimal as a new object. It needs the finite
discipline of comparison: shrink the permitted slip, read the leading
residue, and ask whether that residue can be forced to vanish for every
allowed test. If it can, the path is stationary.

This is why residual vanishing is the grammatical form of stationarity. A
path is not stationary because it looks smooth, short, elegant, or
physically plausible. It is stationary because the grammar has exhausted the
first readable imbalance. Every permitted local slip is tested, and no
linear residue remains. The stationary claim therefore has the shape:
within this gauge, with these endpoints, under these admissible variations,
the first residual vanishes.

A closed flat path gives the simplest picture. Suppose the path returns to
its boundary commitments without accumulating a readable slip. Each local
comparison is flat relative to the gauge, and the closing comparison
restores the initial commitment rather than adding a new debt. The grammar
then reads the whole loop as reducible: not because nothing happened, but
because every permitted difference has already been balanced by the
gauge. The residual ledger has no first-order entry left to spend.

Interpolation gives a useful illustration, provided it remains an
illustration. A thermometer scale may be fixed by two reproducible events
and then supplied with intermediate marks. The intermediate marks are not
new facts about the water. They are readings permitted by the rule that
relates the anchors. If a proposed curve between the anchors introduces a
distinguishable kink that refinement can detect, the kink appears as
residue. If the permitted completion removes every such first-order kink,
the completion reads as stationary relative to the chosen gauge.

The same discipline appears in a calibration path. A sequence of readings
can drift, but not every drift is an admissible variation. Some changes
only rename the scale; some change the relation between anchor events; some
expose a residual imbalance. The gauge identifies the harmless relabelings,
forbids the changes that abandon the committed anchors, and carries the
residue of the remaining changes. Stationarity is the moment when that
carried residue vanishes for all allowed slips.

This first section gives the chapter its central grammar. The finite gauge
permits paths, actions, and variations; it forbids unbounded escape from
the record; it identifies presentations that do not alter the reading; it
carries residuals until they vanish. Once this reader is available, the
next question is forced. If vanishing first residue is stationarity, what
local balance is being read when the residue vanishes?
